\section{The KKT-controlled instance family}\label{sec:family}

\subsection{Dual formulation}
Following~\cite{malick2004dual,qisun2006quadratically}, the Lagrangian dual of~\eqref{eq:ncm}
is the unconstrained problem
\begin{equation}\label{eq:dual}
  \min_{y\in\mathbb{R}^n}\ \theta(y) \coloneqq \tfrac12\fro{\Proj(G+\Diag(y))}^2 - e^{\mathsf T}y ,
  \qquad
  \nabla\theta(y) = \diag\bigl(\Proj(G+\Diag(y))\bigr) - e .
\end{equation}
Here $\Proj$ is the projection onto the positive semidefinite cone. The
gradient is globally Lipschitz with constant 1. If $\ys$ satisfies
$\nabla\theta(\ys)=0$, then $\Xs=\Proj(G+\Diag(\ys))$ is the unique solution
of~\eqref{eq:ncm}. Each evaluation of $\theta$ or $\nabla\theta$ requires one
symmetric eigendecomposition (\EVD) of $C(y)\coloneqq G+\Diag(y)$.

\subsection{Construction and exactness}\label{sec:construction}
The idea is to choose the solution first and then build a $G$ for which that
solution satisfies the KKT conditions exactly. This is the standard technique
of generating instances from prescribed optimality
conditions~\cite{lenard1984randomly,calamai1993new}, and in particular Wei and
Wolkowicz's generator of semidefinite programs with prescribed complementarity
nullity~\cite{wei2010generating}, specialized here to the NCM problem. It also
continues the long tradition of test problems with closed-form solutions
surveyed in Higham's gallery of test matrices~\cite[Ch.~26]{higham2002accuracy}. What is
specific to NCM is the exact projection identity of \cref{prop:exact}, the
control of the rank, multiplicity and separation in \cref{sec:spectrum}, and
the entrywise conditions of \cref{sec:feasibility}.
Their construction does not cover the present problem: it generates the data
of a \emph{linear}-objective SDP and is free to choose the constraint operator
itself (their Algorithm~2.3, Steps 3--4), whereas the NCM problem fixes both
the quadratic objective and the operator $\mathcal{A}=\diag$ with $b=e$, leaving
$G$ as the only free data. The shared element is the idea of prescribing an
optimal primal--dual pair and controlling the dimension of its common
nullspace, their complementarity nullity, which corresponds to $|\beta|$ here.

\begin{definition}[KKT-controlled instance]\label{def:family}
Let $n>r\ge1$. Draw $U\in\mathbb{R}^{n\times r}$ with independent
standard-normal entries and rescale every row to unit Euclidean norm. Set
$\Xs=UU^{\mathsf T}$, so that $\diag(\Xs)=e$, $\Xs\succeq0$ and $\rank\Xs=r$
almost surely. Let $Q_0\in\mathbb{R}^{n\times(n-r)}$ have orthonormal columns
spanning $\ker\Xs$. For $\mu\in\mathbb{R}^{n-r}_{\ge0}$ define
\begin{equation}\label{eq:construction}
  W = Q_0\Diag(\mu)Q_0^{\mathsf T},\qquad
  G = \Xs - W + \Diag(\diag W),\qquad
  \ys = -\diag(W).
\end{equation}
\end{definition}

\begin{proposition}[Exactness]\label{prop:exact}
For every $\mu\ge0$, $\Proj(G+\Diag(\ys))=\Xs$ and $\nabla\theta(\ys)=0$.
Hence $\Xs$ is the unique nearest correlation matrix to $G$.
\end{proposition}
\begin{proof}
By construction, $G+\Diag(\ys)=\Xs-W$. Both $\Xs$ and $W$ are symmetric
positive semidefinite, and $\Xs W=\Xs Q_0\Diag(\mu)Q_0^{\mathsf T}=0$ because
the columns of $Q_0$ lie in $\ker\Xs$. The ranges of $\Xs$ and $W$ are
therefore orthogonal. They need not be complementary: if some $\mu_j=0$, then
$\range W$ is a proper subspace of $\ker\Xs$. Consequently $\Xs-W$ has the
spectral decomposition obtained by combining those of $\Xs$ and $-W$ on
orthogonal subspaces. Its positive part is $\Xs$, and its nonpositive part is
$-W$. Thus $\Proj(\Xs-W)=\Xs$, and
$\nabla\theta(\ys)=\diag(\Xs)-e=0$. Since $\theta$ is convex and
differentiable, $\ys$ minimizes~\eqref{eq:dual}. Uniqueness of the primal
solution is classical~\cite{higham2002nearest,qisun2006quadratically}.
\end{proof}

In floating point, the generator recomputes the projection rather than relying
on the identity. Over the 270 released instances at $n=100$, the KKT residual
is at most $1.7\times10^{-14}$ and the exactness error
$\fro{\Proj(C(\ys))-\Xs}$ at most $8.4\times10^{-14}$.

\subsection{Prescribed spectral structure}\label{sec:spectrum}
Because the ranges are orthogonal, the spectrum of $C(\ys)=\Xs-W$ is the
multiset union
\begin{equation}
  \operatorname{spec}C(\ys) = \{\lambda_1(\Xs),\dots,\lambda_r(\Xs)\}\ \cup\ \{-\mu_1,\dots,-\mu_{n-r}\}.
\end{equation}
In the notation of~\cite{qisun2006quadratically}, let $\alpha,\beta,\gamma$ index the
positive, zero and negative eigenvalues of $C(\ys)$. Then
\begin{itemize}
  \item $|\alpha|=r$, so the rank of the solution is prescribed;
  \item $|\beta|=\#\{j:\mu_j=0\}$ exactly, so the multiplicity of zero
    eigenvalues at the solution is prescribed;
  \item an entry $\mu_j=\delta>0$ produces an eigenvalue $-\delta$, which moves
    toward the nonsmooth boundary of $\Proj$ as $\delta\to0$.
\end{itemize}
With the layout used below ($m$ entries of $\mu$ equal to $\delta$, the rest
equal to $\mu_{\mathrm{bulk}}>0$), the multiplicity and the separation are
therefore not varied independently in their effect on strict
complementarity: $|\beta|=m$ when $\delta=0$ and $|\beta|=0$ otherwise.

\paragraph{Terminology.}
Qi and Sun~\cite{qisun2006quadratically} showed that constraint
nondegeneracy always holds for the NCM problem, so the family does not, and
could not, vary constraint degeneracy. What it varies is \emph{failure of
strict complementarity}: zero eigenvalues of $C(\ys)$ ($|\beta|>0$), and
near-zero eigenvalues (small $\delta$) that approach it. Throughout this
paper, a \emph{degeneracy cell} means one $(r,m,\delta)$ combination in this
sense, and the \emph{strict-complementarity defect} means $|\beta|$.
The generator uses
$\mu=(\mu_{\mathrm{bulk}},\dots,\mu_{\mathrm{bulk}},\delta,\dots,\delta)$ with
$m$ copies of $\delta$ and $\mu_{\mathrm{bulk}}=0.1$. The separation parameter
is $\sep=\delta/\mu_{\mathrm{bulk}}$ (we avoid $\rho$, which denotes the
line-search backtracking factor in \cref{sec:traps,app:algorithms}). The released grid is
\[
  r\in\{5,20,50\},\quad m\in\{1,5,20\},\quad
  \delta\in\{10^{-2},10^{-4},10^{-6},10^{-8},10^{-10},0\},
\]
with five paired seeds per rank. This gives $3\times3\times6=54$ cells and
270 instances at $n=100$. In plots, $\sep$ is shown on a $\log_{10}$ axis and
the case $\sep=0$ is shown separately.

\subsection{Correlation-like inputs: sufficient conditions for $|G_{ij}|\le1$}\label{sec:feasibility}
\Cref{prop:exact} places no restriction on the entries of $G$. For the
instances to resemble invalid correlation matrices, which is the setting in
which NCM problems arise, the generator additionally requires
$\diag(G)=e$ (true by construction), $\lambda_{\min}(G)<0$, and
$|G_{ij}|\le1$. The first two are checked directly. The results below explain
when the third can be expected to hold. Let $P_r$ be the orthogonal projector
onto $\range U$, so that $Q_0Q_0^{\mathsf T}=I-P_r$, and define
\[
  a \coloneqq \max_{i\ne j}|\Xs_{ij}|,\qquad b\coloneqq \max_{i\ne j}|(P_r)_{ij}|.
\]

\begin{proposition}[Deterministic bound, uniform $\mu$]\label{prop:det}
If $\mu=c\,e$ with $c>0$, then $|G_{ij}|\le a+cb$ for all $i\neq j$. In
particular, if $a<1$, then $0<c\le(1-a)/b$ is sufficient for $|G_{ij}|\le1$.
\end{proposition}

\begin{proposition}[Fixed-rank saturation]\label{prop:sat}
Fix $r\ge2$ and let the rows $u_i$ of $U$ be independent and uniform on
$S^{r-1}$. Then $a\to1$ as $n\to\infty$, along every realization.
\end{proposition}

\Cref{prop:sat} says only that $a$ saturates. It does \emph{not} imply that a
fixed $c>0$ eventually violates \cref{prop:det}: at fixed $r$ the quantity $b$
is observed to decrease with $n$ as well, but we do not prove
this, so whether $a+cb$ exceeds 1 depends on the relative
rates of $1-a$ and $b$, which we do not establish. We therefore make no
asymptotic feasibility claim, and the generator relies on entrywise
screening.

\begin{proposition}[High-dimensional regime]\label{prop:hd}
Let the rows of $U$ be independent and uniform on $S^{r-1}$ and let
$\eta\in(0,1)$ be a failure probability (distinct from the near-zero value
$\delta$ in $\mu$). Then with probability at least
$1-\eta$,
\[
  a\;\le\;\sqrt{\frac{2}{r}\Bigl(2\log n+\log\frac{2}{\eta}\Bigr)} .
\]
In particular $a\to0$ in probability whenever $r/\log n\to\infty$.
\end{proposition}

The bound follows from the explicit cap estimate
of~\cite[Lem.~2.2]{ball1997elementary} (subgaussianity in the sense
of~\cite[Thm.~3.4.5]{vershynin2026high} gives the same tail without a
specific constant) and a union bound; see
\cref{app:family}. We state it only for $a$, the quantity that governs
\cref{prop:det} in the regime of interest. A companion bound on $b$ would
require a model for $\range U$ that we do not need here, since the released
family lies in the fixed-rank regime of \cref{prop:sat}, where
\cref{prop:hd} does not apply.

Proofs are given in \cref{app:family}. These results give sufficient
conditions only. Failure of the bound in \cref{prop:det} does not show that
some $|G_{ij}|$ exceeds 1, and for that reason the generator checks every
entry directly instead of relying on the bound. \Cref{prop:det} assumes
uniform $\mu$, whereas the released instances use a mixed $\mu$ (bulk value
$0.1$ together with $m$ copies of $\delta$).
The extension to mixed $\mu$ is left open: the extra off-diagonal term it
introduces depends on the choice of orthonormal basis for $\ker\Xs$ and is not
determined by $(n,r,m,\delta)$ alone (\cref{app:family}). The released
mixed-$\mu$ instances are therefore screened entrywise rather than certified
by a bound.
None of these results says anything about indefiniteness of $G$, which is
checked separately.

\paragraph{Which regime the released family is in.}
\Cref{tab:ab} reports empirical values of $(a,b)$ at $n=100$ over 30 seeds per
rank. The family is in the regime of \cref{prop:sat}, not \cref{prop:hd}. The
rank is fixed and small relative to $n$, $a$ is close to 1 at $r=5$, and the
bound on $c$ implied by \cref{prop:det} shrinks sharply as the rank decreases.

\begin{table}[t]
\centering
\caption{Empirical $(a,b)$ at $n=100$, 30 seeds per rank. The last column is
the sufficient bound $(1-a)/b$ of \cref{prop:det} evaluated at the mean values
of $a$ and $b$.}
\label{tab:ab}
\small
\begin{tabular}{rrrrrr}
\toprule
$r$ & mean $a$ & max $a$ & mean $b$ & max $b$ & $(1-\bar a)/\bar b$\\
\midrule
5  & 0.9905 & 0.9981 & 0.0590 & 0.0708 & 0.162\\
20 & 0.7495 & 0.8340 & 0.1483 & 0.1871 & 1.690\\
50 & 0.5081 & 0.5719 & 0.1856 & 0.2083 & 2.650\\
\bottomrule
\end{tabular}
\end{table}

The table is consistent with the screening outcomes. At
$\mu_{\mathrm{bulk}}=0.1$, the sufficient bound at the mean values at $r=5$
($0.162$) exceeds $0.1$, but at the worst seeds (max $a=0.9981$) the bound
itself falls below $0.1$. In the
paired design (\cref{sec:paired}), 5 of the 7 candidate seeds were accepted at
$r=5$; the 2 rejections were for $|G_{ij}|>1$. At $r=20$ and $r=50$, 5 of 5
were accepted. With $\mu_{\mathrm{bulk}}=1$, the sufficient bound at the mean
values still holds at $r=20$ and $r=50$ ($1.690$ and $2.650$ exceed 1) but
fails at $r=5$, and in an earlier design every $r=5$ candidate was rejected.
Since one bulk level is used at every rank, this is why the generator uses
$\mu_{\mathrm{bulk}}=0.1$.

\subsection{Paired-seed design and screening}\label{sec:paired}
Each candidate seed is screened once per rank against its entire
$(m,\delta)$ block. A seed is accepted only if all $3\times6$ instances built
from it pass the checks ($\diag(G)=e$, $|G_{ij}|\le1$, $\lambda_{\min}(G)<0$,
and a recomputed KKT residual and exactness error each below $10^{-9}$). The
accepted seeds are then reused across every $(m,\delta)$ cell. Varying the
separation therefore changes only $\mu$, with $U$ held fixed, so comparisons
along $\delta$ are within-seed. Rejections are counted and reported, not
silently discarded.

Because of the screening, the $r=5$ sample is conditioned differently from the
$r=20$ and $r=50$ samples: at $r=5$ it excludes the seeds whose instances have
an off-diagonal entry above 1 (at $n=100$, both rejections occurred at $m=20$,
$\delta=10^{-2}$). These are not simply the seeds with the largest $a$, since
the rejection depends on the basis-dependent kernel term as well.
Comparisons across ranks should be read with this in mind.

\paragraph{Extent of the family.}
Every synthetic instance is close to feasible: $\lambda_{\min}(G)$ lies between
$-0.075$ and $-0.017$ at $n=100$ and between $-0.047$ and $-0.011$ at $n=500$.
A larger bulk level would give more strongly invalid inputs, but it fails the
entrywise screen at $r=5$, and the family uses a single bulk level at every
rank; the degree of invalidity is therefore not a controlled factor here.
Findings that depend on it (\cref{sec:thesis,sec:us2105}) rest on perturbed
matrices whose reference is computed, not exact.

\paragraph{Regeneration.}
The seed of candidate $k$ at rank $r$ and dimension $n$ is $100000n+1000r+k$.
The multiplicity is limited by the kernel dimension, $m\le n-r$, and the $m$
entries equal to $\delta$ are attached to the last $m$ columns of $Q_0$, the
kernel block of the eigenvector basis of $\Xs$ returned by the symmetric
eigensolver, which is the point at which a
different LAPACK build can yield a different member of the family
(\cref{app:artifact}).

\subsection{Release}
We release the 270 paired instances ($G$, $\Xs$, $\ys$, SHA-256 hashes of the
upper triangle of $G$, and per-instance screening diagnostics), the generator,
and the seed-screening log. \Cref{app:artifact} lists the files.
We also release a paired set at $n=500$, generated on the same grid. All 270
instances passed screening; seed acceptance was 5 of 7 candidates at $r=5$
(both rejections for $|G_{ij}|>1$) and 5 of 5 at $r=20$ and $r=50$, with
$\max_{i\ne j}|G_{ij}|=0.99987$ among the accepted $r=5$ instances. The
ranking study of \cref{sec:n500} uses all five paired seeds per rank.
